% !TeX root = ../main.tex
\section{Language and Determinization}
\label{sec:language}

This section introduces the probabilistic language on which our
determinization transformation operates. We first present its syntax and operational semantics, and then define
determinization as a source-to-source transformation over annotated expressions. Finally, we define the output measure that characterizes the distribution of successfully returned values.

\paragraph{Syntax.}
\Cref{fig:syntax} presents the syntax of our language, a call-by-value
functional language with recursion, pairs, sums, lists, and probability distributions. The expression $\kw{reject}$ represents a failed execution path used to encode conditioning on observations.

\todo[inline]{TODO: explain why reject and diverge are treated similarly??}

Each distribution is annotated with one of three modes:
$\G$, $\E$, or $\M$. Both $\G$ and $\E$ draw samples, but $\G$ marks sampling sites that remain stochastic under determinization, whereas $\E$ marks those eligible to be replaced by their means. In particular, the type system in \cref{sec:type-system} determines which sampling sites may be assigned mode $\E$. The mode $\M$, used in the output of determinization, evaluates a distribution expression to its mean.

% !TeX root = ../main.tex
\begin{figure}[H]
\centering
\small
\begin{align*}
 e &::= x\mid r\;(r\in\Real)\mid ()\mid\kw{true}\mid\kw{false}\mid\kw{reject}\\
   &\quad\mid \lambda x.e\mid\recfn f x e\mid e\,e
       \mid\letin x e e\\
   &\quad\mid \ite e e e\mid -e\mid e+e\mid e\cdot e\mid e/e\mid e<e\\
   &\quad\mid \langle e,e\rangle\mid\kw{fst}\,e\mid\kw{snd}\,e\\
   &\quad\mid \kw{inl}\,e\mid\kw{inr}\,e\mid\scase e x{e_1}y{e_2}\\
   &\quad\mid []\mid e::e\mid\lcase e{e_0}x{xs}{e_1}\\
   &\quad\mid \draw{uniform}\alpha{e,e}\mid\draw{gaussian}\alpha{e,e}\\
   &\quad\mid \draw{poisson}\alpha e\mid\draw{exponential}\alpha e\\
   &\quad\mid \draw{bernoulli}\alpha e\mid\draw{discrete}\alpha e\\
   &\quad\mid \draw{beta}\alpha{e,e}\mid\draw{gamma}\alpha{e,e}.\\[2pt]
 \alpha &::= \G\mid\E\mid\M
\end{align*}
\caption{Language grammar. $\G$ and $\E$ denote sampling; $\M$ denotes mean evaluation.}
\label{fig:syntax}
\Description{Expressions include literals, functions and recursion, arithmetic, conditionals, products, sums, lists, rejection, and eight distribution primitives. The G, E, and M annotations are listed below the expressions.}
\end{figure}


\paragraph{Operational semantics.}
\Cref{fig:semantics} gives the call-by-value operational semantics.
Values and evaluation contexts in
\cref{fig:semantics}(\subref{fig:semantics-values}) specify which
expressions have finished evaluating and which subexpression is evaluated
next. In particular, a distribution's arguments are evaluated to values
before it is sampled or its mean is computed.

The deterministic reduction rules in
\cref{fig:semantics}(\subref{fig:semantics-reduction}) handle the standard
functional constructs by substitution and case analysis. We omit the
routine arithmetic and comparison rules. Distribution expressions in
mode $\M$, introduced by determinization, also reduce deterministically
to their means:
\[
  D^\M(\bar v)
  \rightsquigarrow
  \operatorname{Mean}_D(\bar v)
  \qquad
  \text{when }\operatorname{Domain}_D(\bar v),
\]
where the primitive distribution table in
\cref{fig:semantics}(\subref{fig:semantics-distributions}) specifies
$\operatorname{Domain}_D$ and $\operatorname{Mean}_D$ for each supported
distribution. A distribution application outside its domain becomes
stuck, as does division by zero. In $\kw{gaussian}(u,v)$, $v$ denotes
the variance; the parameters $\lambda$ of the exponential and gamma
distributions are rates. The distribution
$\kw{discrete}([p_i]_{i<n})$ returns $i$ with probability $p_i$ for
$0\leq i<n$, and returns $n$ with the remaining probability
$1-\sum_{i<n}p_i$.

In contrast to mode $\M$, modes $\G$ and $\E$ both sample from
$D(\bar v)$ when the arguments satisfy the domain condition $\operatorname{Domain}_D(\bar v)$.
Since sampling is stochastic, neither mode has a deterministic reduction
rule. Their identical sampling behavior is instead specified by the
output-measure rules in
\cref{fig:semantics}(\subref{fig:semantics-output}), which we describe
later in this section.

% !TeX root = ../main.tex
\begin{figure}[p]
\centering
\small
\setlength{\jot}{2pt}
\newcommand{\semanticsHeading}[2]{%
  \par\nointerlineskip\phantomsubcaption\label{#2}%
  \textbf{(\thesubfigure) #1}\par\nobreak\vspace{6pt}\nointerlineskip}
\begin{subcaptiongroup}
\semanticsHeading{Values\leanref{values} and evaluation contexts\leanref{evaluation-contexts}}{fig:semantics-values}
\(\displaystyle
\begin{aligned}
 v&::=r\mid()\mid\kw{true}\mid\kw{false}\mid\lambda x.e\mid\recfn f x e
       \mid\langle v,v\rangle\mid\kw{inl}\,v\mid\kw{inr}\,v\mid[]\mid v::v\\
 C&::=[\,]\mid C\,e\mid v\,C\mid\letin x C e\mid\ite C{e_1}{e_2}\\
  &\quad\mid\langle C,e\rangle\mid\langle v,C\rangle\mid\kw{fst}\,C\mid\kw{snd}\,C
       \mid\kw{inl}\,C\mid\kw{inr}\,C\\
  &\quad\mid C::e\mid v::C\mid\scase C x{e_1}y{e_2}\mid\lcase C{e_0}x{xs}{e_1}\\
  &\quad\mid -C\mid C\mathbin\circ e\mid v\mathbin\circ C
       \mid\draw D\alpha{\bar v,C,\bar e},
       \qquad \circ\in\{+,\cdot,/,<\},\quad \alpha\in\{\G,\E,\M\}.
\end{aligned}
\)\par
\par\vspace{11pt}
\semanticsHeading{Deterministic reduction\leanref{reduction}}{fig:semantics-reduction}
\begingroup
\setlength{\jot}{0pt}
\(\displaystyle
\begin{aligned}[t]
 (\lambda x.e)\,v&\rightsquigarrow e[x\mapsto v]\\
 (\recfn f x e)\,v&\rightsquigarrow e[x\mapsto v,f\mapsto\recfn f x e]\\
 \letin x v e&\rightsquigarrow e[x\mapsto v]\\
 \kw{reject}&\rightsquigarrow\kw{reject}\\
 \ite{\kw{true}}{e_1}{e_2}&\rightsquigarrow e_1\\
 \ite{\kw{false}}{e_1}{e_2}&\rightsquigarrow e_2\\
 \kw{fst}\,\langle v_1,v_2\rangle&\rightsquigarrow v_1\\
 \kw{snd}\,\langle v_1,v_2\rangle&\rightsquigarrow v_2\\[3pt]
 \scase{\kw{inl}\,v}x{e_1}y{e_2}&\rightsquigarrow e_1[x\mapsto v]\\
 \scase{\kw{inr}\,v}x{e_1}y{e_2}&\rightsquigarrow e_2[y\mapsto v]\\
 \lcase{[]}{e_0}x{xs}{e_1}&\rightsquigarrow e_0\\
 \lcase{v::w}{e_0}x{xs}{e_1}&\rightsquigarrow e_1[x\mapsto v,xs\mapsto w]\\[3pt]
 \avg D{\bar v}&\rightsquigarrow \operatorname{Mean}_D(\bar v)\quad(\operatorname{Domain}_D(\bar v)).
\end{aligned}
\)\par
\endgroup
\par\vspace{11pt}
\semanticsHeading{Primitive distributions\leanref{primitive-distributions}}{fig:semantics-distributions}
\begingroup
\renewcommand{\arraystretch}{1.08}
\setlength{\heavyrulewidth}{0.5pt}
\begin{tabular*}{0.72\linewidth}{@{\extracolsep{\fill}}lll@{}}
\toprule
Distribution $D(\bar v)$ & $\operatorname{Domain}_D(\bar v)$ & $\operatorname{Mean}_D(\bar v)$\\
\midrule
$\kw{uniform}(a,b)$ & $a\le b$ & $(a+b)/2$\\
$\kw{gaussian}(u,v)$ & $v\ge0$ & $u$\\
$\kw{poisson}(\lambda)$ & $\lambda\ge0$ & $\lambda$\\
$\kw{exponential}(\lambda)$ & $\lambda>0$ & $1/\lambda$\\
$\kw{bernoulli}(p)$ & $0\le p\le1$ & $p$\\
$\kw{discrete}([p_i]_{i<n})$ & $p_i\ge0,\ \sum_{i<n}p_i\le1$ & $n+\sum_{i<n}(i-n)p_i$\\
$\kw{beta}(a,b)$ & $a>0,\ b>0$ & $a/(a+b)$\\
$\kw{gamma}(k,\lambda)$ & $k>0,\ \lambda>0$ & $k/\lambda$\\
\bottomrule
\end{tabular*}
\endgroup
\par\vspace{11pt}
\semanticsHeading{Output measures\leanref{output-measures}}{fig:semantics-output}
\begingroup
\settowidth{\dimen0}{$\displaystyle\int_{\Real}\out{C[r]}n(B)\,D(\bar v)(dr)$}
\edef\outputTermWidth{\the\dimen0}
\settowidth{\dimen2}{$e=C[\draw D m{\bar v}],\enspace\operatorname{Domain}_D(\bar v),\enspace m\in\{\G,\E\}$}
\edef\outputConditionWidth{\the\dimen2}
\setlength{\jot}{0pt}
\(\displaystyle
\begin{aligned}
 \out e0(B)&=
 \left\{\begin{array}{@{}w{l}{\outputTermWidth}@{\quad}w{l}{\outputConditionWidth}@{\quad}l@{}}
  \delta_r(B) & e=r,\quad r\in\Real & \text{(Return real numbers)}\\
  0 & \text{otherwise} & \text{(No real result yet)}
 \end{array}\right.\\[5pt]
 \out e{n+1}(B)&=
 \left\{\begin{array}{@{}w{l}{\outputTermWidth}@{\quad}w{l}{\outputConditionWidth}@{\quad}l@{}}
  \out{C[e_2]}n(B) & e=C[e_1],\quad e_1\rightsquigarrow e_2
    & \text{(Deterministic step)}\\[5pt]
  \displaystyle\int_{\Real}\out{C[r]}n(B)\,D(\bar v)(dr)
    & e=C[\draw D m{\bar v}],\enspace\operatorname{Domain}_D(\bar v),\enspace m\in\{\G,\E\}
    & \text{(Average over samples)}\\[7pt]
  0 & \text{otherwise} & \text{(No output at this depth)}
 \end{array}\right.\\[5pt]
 \law e(B)&=\sum_{n=0}^{\infty}\out e n(B).
\end{aligned}
\)\par
\endgroup
\end{subcaptiongroup}
\captionsetup{skip=10pt}
\caption{Operational semantics. In (b), we omit the usual arithmetic and comparison reductions; division by zero gets stuck. In (d), the output measures are subprobability measures: executions that do not successfully return a real value contribute no mass. Moreover, deterministic reductions proceed according to $\rightsquigarrow$,
while $\G$- and $\E$-mode sampling steps integrate over
all possible sampled values.}
\label{fig:semantics}
\Description{The semantics consists of values and evaluation contexts, deterministic reductions, primitive distributions with domains and means, and output measures summed over exact termination depths.}
\end{figure}


\paragraph{Determinization} We now formally define our determinization function below.

\begin{definition}[Determinization\leanref{determinization}]
\label{def:determinization}
Let $\mathcal{A}=\{\G,\E,\M\}$ be the set of distribution modes. The
determinization transformation $\trans{\cdot}$ is defined structurally on
expressions, replacing each $\E$-mode distribution by its $\M$-mode
counterpart and leaving all other constructs unchanged. In particular, for
every distribution $D$, mode $\alpha\in\mathcal{A}$, and argument
expressions $\bar e=(e_1,\ldots,e_n)$,
\[
\trans{D^\alpha(\bar e)}
=
\begin{cases}
D^\M(\trans{\bar e}) & \text{if } \alpha=\E,\\
D^\alpha(\trans{\bar e}) & \text{if } \alpha\in\{\G,\M\},
\end{cases}
\]
where
\[
\trans{\bar e}
=
(\trans{e_1},\ldots,\trans{e_n}).
\]
For every non-distribution constructor $C$,
\[
\trans{C(e_1,\ldots,e_n)}
=
C(\trans{e_1},\ldots,\trans{e_n}),
\]
and atomic expressions are unchanged (see~\Cref{sec:appendix-determinization}).
\end{definition}

Determinization replaces $\E$-mode sampling with mean evaluation and recursively transforms the distribution's arguments, preserving their call-by-value evaluation order. When all sampling sites are annotated with $\E$, determinization eliminates all randomness, yielding a deterministic program. We refer to this as \emph{full determinization}. When both $\E$- and $\G$-mode sampling sites occur, determinization eliminates the former while retaining the latter. Accordingly, we refer to this as \emph{partial determinization}. In this case, the resulting program may be discrete or remain continuous, depending on the randomness retained. In particular, if all retained sampling sites are discrete, the transformed program is discrete, even if the original program contains continuous distributions. Conversely, retaining continuous sampling sites may leave the transformed program continuous. Let us consider several simple examples below to illustrate determinization in action.

\begin{example}
In this example, determinization first replaces the outer $\E$-annotated Gaussian by its $\M$-annotated counterpart, and then recursively transforms its arguments (trivial recursive steps on constants, which are unchanged by determinization, are omitted for brevity). The nested uniform distribution is annotated with $\E$, so it also gets replaced by its $\M$-annotated counterpart:
\begin{align*}
\trans{
  \kw{gaussian}^{\E}
    \bigl(\kw{uniform}^{\E}(0,1),1\bigr)}
&=
\kw{gaussian}^{\M}
  \bigl(
    \trans{\kw{uniform}^{\E}(0,1)},
    1
  \bigr)
\\
&=
\kw{gaussian}^{\M}
  \bigl(
    \kw{uniform}^{\M}(0,1),
    1
  \bigr).
\end{align*}
Thus, the Gaussian and uniform draws are determinized to their means.
\end{example}

\begin{example}
Consider the following partial determinization example, where the draws are dependent. The $\G$-annotated uniform draw for $x$ remains stochastic, while the conditional Gaussian draw for $y$ is determinized. Importantly, its mean $2x+1$ may still depend on the sampled value of $x$:
\[
\trans{\left(
\begin{aligned}
&\kw{let}\ x=\kw{uniform}^{\G}(0,1)\ \kw{in}\\
&\kw{let}\ y=\kw{gaussian}^{\E}(2x+1,\sigma^2)\ \kw{in}\\
&x*y
\end{aligned}
\right)}
=
\begin{aligned}
&\kw{let}\ x=\kw{uniform}^{\G}(0,1)\ \kw{in}\\
&\kw{let}\ y=\kw{gaussian}^{\M}(2x+1,\sigma^2)\ \kw{in}\\
&x*y.
\end{aligned}
\]
Thus, determinization removes the randomness introduced by $y$ without removing the randomness on which its mean depends.
\end{example}

We note that replacing a sample by its mean is not valid in every context. For example,
if $X$ is uniformly distributed on $[0,1]$, then $\mathbb{E}[X^2]=\frac{1}{3}$ but $\mathbb{E}[X]^2=\frac{1}{4}$.

The type system therefore prevents the sample in
\[
  \kw{let}\ x=\kw{uniform}^{\E}(0,1)\ \kw{in}\ x*x
\]
from being assigned mode $\E$. We present these restrictions in
\cref{sec:type-system}. Thus, the transformation defined above is
syntactic: it performs the requested replacements, while the type system
establishes that those replacements are semantically justified.

\paragraph{Output measures.}
As described above, the reduction rules describe how an individual execution proceeds. To
describe the behavior of all possible executions, we lift these rules to a
measure over returned real values, as defined in \cref{fig:semantics}(\subref{fig:semantics-output}). For a closed expression $e$ and
$n\in\mathbb{N}$, $\out e n$ is a measure on the standard Borel space
$\mathbb{R}$. For every Borel-measurable set $B\subseteq\mathbb{R}$,
$\out e n(B)$ is the probability that $e$ returns a value in $B$ after
exactly $n$ reduction steps. Deterministic steps continue with the resulting
expression, while sampling steps integrate over every possible sampled
value. Summing over all finite return depths gives the complete output law of $e$.

The output law may have total mass less than one: for example, executions that do not
return a real number, contribute no mass. Consequently, $\law e(\mathbb{R})$
is the probability that $e$ successfully returns a real value. 
This semantics provides the basis for the results that follow. The output measures give us the mathematical object used in \cref{sec:soundness} to prove that determinizing a
well-typed program preserves its expected returned value while reducing
variance under the stated assumptions.

\todo[inline]{TODO: fix and explain domain safety. In:

Operational semantics: explain when an operation becomes stuck

Output measures: explain that stuck executions contribute no output mass.

Domain safety: define it as requiring operations to remain within their domains almost surely at every finite stage of execution. Clarify that probability-zero failures are allowed, while positive-probability failures are not.}
\begin{remark}[Domain safety\leanref{domain-safety}]
We assume programs are \emph{domain-safe}: almost surely, every executed
distribution has parameters in its domain and every executed division has a
nonzero divisor.  Thus execution cannot become stuck because a primitive
operation is undefined.  We show below that this assumption is necessary.
\end{remark}

\paragraph{Why domain safety is necessary.}
The preceding guarantees rely essentially on domain safety. Without it, replacing an $\E$-mode sample by its mean can change whether a later operation is defined, and therefore change which executions contribute to the return-conditioned distribution.

For example,
\begin{lstlisting}
let u = gaussian[E](0, 1) in
uniform[E](0, u)
\end{lstlisting}
is not domain-safe: executions with $u<0$ become stuck. Conditioning on successful return therefore implicitly conditions $u$ to be positive, giving
\[
\Expect_{\mathrm{ret}}[e]
=
\frac{1}{2}\Expect[u\mid u>0]
=
\frac{1}{\sqrt{2\pi}}.
\]
Determinization instead replaces $u$ by its unconditional mean $0$, yielding
\[
\Expect_{\mathrm{ret}}[\trans e]=0.
\]
Thus domain safety prevents determinization from changing which executions are semantically valid and, consequently, which executions contribute to the return-conditioned output distribution.
